\chapter{Prove Podman Works}
\label{chap:prove-podman-works}

Before building anything, establish the current world state.

A common mistake is to begin changing a system before understanding how it is configured. When something later breaks, there is no reliable record of what changed.

This chapter has a simple goal:

\begin{quote}
Ask the computer what it knows before asking it to do work.
\end{quote}

\section{Learning goals}

\begin{itemize}
  \item Verify that Podman is installed.
  \item Observe how Podman is configured.
  \item Distinguish between the client and container host.
  \item Collect evidence before making changes.
  \item Create durable troubleshooting artifacts.
\end{itemize}

\section{Observe before changing}

For this chapter we run a series of discovery commands.

\begin{lstlisting}[language=bash]
podman --version

podman machine list

podman system connection list

wsl -l -v

podman info

podman images

podman ps -a

podman system df
\end{lstlisting}

Do not worry about understanding every field.

Instead, focus on the types of questions each command helps answer.

\section{Why collect evidence?}

The following process works surprisingly well:

\begin{enumerate}
  \item Observe.
  \item Record.
  \item Interpret.
  \item Change one thing.
  \item Observe again.
\end{enumerate}

Many troubleshooting sessions become confusing because multiple changes are made before any evidence is collected.

A written record allows us to compare the system before and after a modification.

\section{Observations from a real workstation}

The following observations were collected from an actual course workstation using the evidence collection script included in this repository.

\subsection{Podman is installed}

The command:

\begin{lstlisting}[language=bash]
podman --version
\end{lstlisting}

reported:

\begin{lstlisting}
podman version 6.1.0
\end{lstlisting}

This confirms that the Podman client is installed and can execute successfully.

\subsection{A Podman machine exists}

The command:

\begin{lstlisting}[language=bash]
podman machine list
\end{lstlisting}

reported a machine named:

\begin{lstlisting}
podman-machine-default
\end{lstlisting}

running under WSL.

This tells us that containers are not running directly inside Windows. A Linux environment is responsible for hosting containers.

\subsection{Windows and Linux cooperate}

The command:

\begin{lstlisting}[language=bash]
podman info
\end{lstlisting}

reported information from both Windows and Linux.

Important observations included:

\begin{itemize}
  \item Windows acts as the client.
  \item Linux acts as the container host.
  \item WSL2 provides the underlying platform.
\end{itemize}

This is a useful reminder that the command prompt you see is not necessarily the operating system executing the container workload.

\subsection{Images already exist}

The workstation already contained several container images including:

\begin{itemize}
  \item Alpine
  \item BusyBox
  \item Debian
\end{itemize}

An image is not a running program.

An image is a template used to create containers.

We will build our own images later in this text.

\subsection{No containers were running}

The command:

\begin{lstlisting}[language=bash]
podman ps -a
\end{lstlisting}

returned no containers.

This is an important observation.

The machine contains images, but no active workloads.

Images and containers are different concepts and should not be confused.

\subsection{Storage is already being used}

The command:

\begin{lstlisting}[language=bash]
podman system df
\end{lstlisting}

reported storage usage associated with container images.

Even when containers are not running, images still occupy disk space.

Containers, images, networks, and volumes are all resources that consume storage and should be managed intentionally.

\section{A useful failure}

During the investigation, several Podman commands initially failed.

The cause was simple:

\begin{itemize}
  \item Podman was installed.
  \item The Podman machine existed.
  \item The Podman machine was not running.
\end{itemize}

After collecting evidence, the machine was started and the investigation continued.

This experience reinforces an important habit:

\begin{quote}
Failures are data.
\end{quote}

An error message is evidence.

Record it before attempting a fix.

Many troubleshooting breakthroughs come from carefully reading the initial failure rather than immediately searching for a workaround.

\section{Keeping durable artifacts}

This repository contains the script:

\begin{lstlisting}[language=bash]
scripts/collect-podman-evidence.ps1
\end{lstlisting}

The script records:

\begin{itemize}
  \item Podman version information.
  \item Podman machine status.
  \item Connection information.
  \item WSL status.
  \item Container images.
  \item Existing containers.
  \item Storage utilization.
\end{itemize}

The output is written to the \texttt{reports} directory.

Treat these reports as evidence.

They help answer the question:

\begin{quote}
What did the system look like before we changed it?
\end{quote}

\section{What we learned}

By the end of this chapter we have not built a single container.

That is intentional.

Before creating anything, we learned how to establish the current state of the environment, collect evidence, and create durable records.

Those habits will be useful long after the specific Podman commands are forgotten.

\checkpoint{Save command output before making changes. Establish the world state first.}